\section{The actual original rank-one chart algebra}

\begin{lemma}[Homogeneous rank-one polynomials and actual dehomogenization]
\label{lem:fc-native-rank-one-homogeneous-polynomial}
\leanok
\lean{EnvelopeProperty.MultiModel.rankOne_homogeneous_eq_monomial, EnvelopeProperty.MultiModel.rankOne_dehom_eq_zero}
\uses{lem:fc-native-model-total-components}
A homogeneous degree-$d$ polynomial in the original single variable is
its original degree-$d$ coefficient times that variable to degree $d$.
Dehomogenization at an original nonzero coordinate is therefore injective
on that homogeneous piece.
\end{lemma}
\begin{proof}
Every original exponent in one variable is the single exponent at its
unique index, with value equal to its total degree. Homogeneity makes
every other coefficient zero. Dehomogenization of the remaining monomial
is its coefficient divided by the original nonzero coordinate to the
same degree, so it vanishes only if the coefficient vanishes.
\end{proof}

\begin{theorem}[Whole native homogeneous chart equals the original chart]
\label{thm:fc-native-rank-one-chart-comparison}
\leanok
\lean{EnvelopeProperty.MultiModel.nativeRankOneChartInput, EnvelopeProperty.MultiModel.nativeRankOneChartInput_coordinate, EnvelopeProperty.MultiModel.nativeRankOneOrdinaryDehom, EnvelopeProperty.MultiModel.nativeRankOneChartDehom, EnvelopeProperty.MultiModel.nativeRankOneChartDehom_mk, EnvelopeProperty.MultiModel.nativeRankOneChartDehom_injective, EnvelopeProperty.MultiModel.nativeRankOneChartDehom_surjective, EnvelopeProperty.MultiModel.nativeRankOneChartEquiv, EnvelopeProperty.MultiModel.nativeRankOneOriginalChartIso, EnvelopeProperty.MultiModel.nativeRankOneChartEquiv_coefficients, EnvelopeProperty.MultiModel.nativeRankOneOriginalChartIso_structure}
\uses{lem:fc-native-rank-one-model-charts, lem:fc-native-rank-one-homogeneous-polynomial, lem:gl-charts}
The actual degree-zero homogeneous localization at an original nonzero
coordinate is isomorphic as a whole ring to the SAME original chart
algebra. The resulting actual native basic-open scheme is the spectrum
of that original chart algebra.
\end{theorem}
\begin{proof}
The original dehomogenization maps the chosen coordinate to one. Its
actual ring homomorphism into the original chart algebra consequently
extends to the ordinary localization, then restricts to its true
degree-zero homogeneous subring. A homogeneous fraction maps to the
original numerator's dehomogenization, since its denominator maps to
one. If the fraction maps to zero, injectivity on its homogeneous
numerator forces that numerator to be zero, so the original fraction
vanishes. Conversely, every original chart element has the original
same-degree numerator/coordinate-denominator presentation; the matching
native homogeneous fraction maps to it. The resulting bijective actual
ring homomorphism gives the whole-ring isomorphism, and contravariant
Spec followed by the genuine native basic-open isomorphism gives the
whole original chart scheme isomorphism.
\end{proof}

The literal original coefficient inclusion is preserved by the whole chart ring isomorphism. The inverse of the actual original chart scheme isomorphism, followed by the actual original coordinate-open inclusion and original structure morphism, is exactly the spectrum map of the SAME original chart coefficient inclusion. This follows from the native Proj homogeneous-chart structure triangle and the proved whole-ring coefficient identity.
